\slajdid{09_3-s003}
\begin{frame}[label=09_3-s003]
\gusttekst\setlength{\abovedisplayskip}{3pt}\setlength{\belowdisplayskip}{3pt}\setlength{\abovedisplayshortskip}{2pt}\setlength{\belowdisplayshortskip}{2pt}Kapacitivnosti MOS FET tranzistora su
\begin{columns}[c,onlytextwidth]\column{.37\textwidth}\centering\begin{tikzpicture}[x=.75cm,y=.75cm,font=\sitneoznake,>=Latex]
\ctikzset{bipoles/length=.65cm}
\draw(1.45,1.6)--(1.45,2.6);\draw(1.3,1.55)--(1.3,2.65);
\draw(1.45,2.55)--(2,2.55)--(2,4.1)node[above]{D};
\draw(1.45,1.65)--(2,1.65)--(2,0)node[below]{S};
\draw(-.8,1.85)node[left]{G}--(1.3,1.85);
\draw[->](3.9,2.1)node[right]{B}--(1.45,2.1);
\draw(0,1.85)to[C](0,3.5)--(2,3.5);
\draw(0,1.85)to[C](0,.55)--(3.25,.55);
\draw(2,3.9)--(3.25,3.9)to[C,l=$C_{DB}$](3.25,2.1)to[C,l=$C_{SB}$](3.25,.55);
\fill(0,1.85)circle(.045);\fill(3.25,2.1)circle(.045);
\draw(.85,1.85)--(.85,.7);\draw[preaction={draw=white,line width=2pt}](.85,.7)arc[start angle=90,end angle=270,radius=.15];\draw(.85,.4)to[C,l_=$C_{GB}$](.85,-.9)--(2.65,-.9)--(2.65,.4);
\draw[preaction={draw=white,line width=2pt}](2.65,.4)arc[start angle=270,end angle=90,radius=.15];\draw(2.65,.7)--(2.65,2.1);
\node[left]at(-.25,2.7){$C_{GD}$};\node[left]at(-.25,1.15){$C_{GS}$};\end{tikzpicture}
\column{.61\textwidth}
$C_{GB}$ - prirodna kapacitivnost između gejta i osnove za MOS FET tranzistor– dovođenje naelektrisanja iz osnove u kanal
\par\medskip
$C_{GD}$, $C_{GS}$ – Oblast gejta mora da pokrije oblast kanala. Međutim teško je obezbediti tehnološki da bude „samo“ iznad kanala. Prekriće delimično i oblasti sorsa i dreja, a time će se pojaviti i ova kapacitivnost.
\par\medskip\centering\begin{tikzpicture}[x=.8cm,y=.65cm,font=\sitneoznake,>=Latex]
\draw(-.5,0)--(4.1,0);\draw(-.5,-1.6)--(4.1,-1.6);\node[right]at(-.4,-1.25){p osnova};
\foreach \x/\d/\l in{.4/p+/B,1.5/n+/S,3.2/n+/D}{\draw(\x-.3,0)--(\x-.3,-.25)arc[start angle=180,end angle=360,x radius=.3,y radius=.35]--(\x+.3,0);\node at(\x,-.25){\d};\draw(\x-.12,0)rectangle(\x+.12,.15);\draw(\x,.15)--(\x,.6)node[above]{\l};}
\draw(1.75,.02)rectangle(2.95,.17);\draw(2.35,.17)--(2.35,.6)node[above]{G};
\node at(2.4,-1.05){preklapanje};\draw[->](2.12,-.83)--(1.78,-.05);\draw[->](2.65,-.83)--(2.91,-.05);
\end{tikzpicture}
\par\medskip\raggedright
$C_{DB}$, $C_{SB}$ – Kapacitivnosti spojeva osnove i drejna (pn spoj – dioda), odnosno osnove i sorsa (pn spoj – dioda)
\end{columns}
\end{frame}
